\section{Discussion and limitations}
\label{sec:disc}

\textbf{What determines the error.} In our experiments the RMSE of the regression GATE estimator is well described by $\sqrt{B^2+\hat\kappa^2L/n}$ (Fig.~\ref{fig:ext}d). $L$ depends only on graph and design and can be computed \emph{before} an experiment (Prop.~\ref{prop:lev}): it grows with the harmonic-mean degree and $1/(p(1-p))$ and is cut by clusterings that keep neighbourhoods together. $B$ is what the regression cannot represent: the response slope away from the exposures the design visits, interference beyond the modelled neighbourhood, and effect heterogeneity correlated with a unit's exposure variance. More data does not cure a nonzero $B$; a larger model helps only where the design or low-degree units provide support, and otherwise trades bias for leverage. Practically: (i) use $L$ to size experiments and pick $p$ (near $\frac12$) and the clustering; (ii) prefer cluster randomization when the graph is clusterable, balancing bias against the number of independent units; (iii) comparing with richer specifications is a natural sensitivity check, but their leverage was up to $20\times$ larger, so only large biases are detectable (not evaluated formally); (iv) intervals inherit $B$.

\textbf{Limitations.} \emph{Synthetic outcomes:} four stylized mechanisms with chosen parameters ($\gamma{=}1$, $\lambda\le0.8$); bias magnitudes belong to these, not to platforms, and real interference may combine them or add attribute heterogeneity and dynamics. \emph{Known graph} except in E7 (one of many mis-measurement models). \emph{Theory:} Props.~\ref{prop:decomp}--\ref{prop:mh} are first-order design-expectation statements without rates; there is no closed form for $B$ under cluster designs, and the leverage formula underpredicts RMSE in 73\% of cluster-design runs (median $+10\%$, up to $3.7\times$). \emph{Estimators:} hand-specified exposure features, two heterogeneity fixes and one boosting model (few replications); we did not implement the covariate selection of~\citep{han2023model}, regularized fits or optimized clusterings (Louvain, unequal sizes). \emph{Inference:} three variance estimators, none correcting bias. \emph{Scale:} $\le1.2\times10^4$ nodes, 6 worlds per setting, bars are min--max not confidence intervals. \emph{Pilot:} \textsc{wls-s}/\textsc{strat-s} were designed after the heterogeneity pilot, so Table~\ref{tab:het} is in-sample for them. Negative results (little clustering benefit on weakly clusterable graphs, no help from quadratic terms, two-step exposure raising RMSE, anticonservative HAC, biased boosting) are reported in Sec.~\ref{sec:results}.
